\chapter{B05 证据矩阵与验收表}\label{sec:b05-appendix}
\section{B05 单元覆盖}

\begin{longtable}{@{}>{\raggedright\arraybackslash}p{0.075\textwidth}>{\raggedright\arraybackslash}p{0.255\textwidth}>{\raggedright\arraybackslash}p{0.32\textwidth}>{\raggedright\arraybackslash}p{0.20\textwidth}@{}}
  \caption{B05 的 11 个内容单元、证据用途与处置。}\label{tab:b05-coverage}\\
  \toprule
  单元 & 主要 TeX 来源 & 进入正文的内容 & 处置与目标 \\
  \midrule
  \endfirsthead
  \toprule
  单元 & 主要 TeX 来源 & 进入正文的内容 & 处置与目标 \\
  \midrule
  \endhead
  \bottomrule
  \endfoot
  U035 & \texttt{构造胶子准算符.tex} & 六振幅、投影、重整化分类、格点流 & integrate；胶子算符、TMD 与自旋 \\
  U041 & \texttt{格点QCD中的TMD\_PDF.tex} & quark TMD、staple、软核、CS、匹配状态 & integrate；TMD 定义、软结构、匹配与验收 \\
  U063 & \texttt{格点QCD中的胶子极化.tex} & 轴反常、螺旋度算符、赝 ITD、系统误差 & integrate；胶子极化、TMD 与自旋 \\
  U064 & \texttt{格点QCD中的胶子算符.tex} & U035 扩展、hybrid/RI-MOM、应用边界 & integrate；胶子算符、TMD 与重正化 \\
  U076 & \texttt{格点上计算胶子准算符.tex} & 实现管道、蒸馏、非连通三点函数 & integrate；TMD、准 TMD 与软结构；本证据附录 \\
  U080 & \texttt{胶子PDF幻灯片.tex} & 因子化与准/赝 PDF 教学顺序 & deduplicate；U083/U086 \\
  U083 & \texttt{胶子PDF连续极限.tex} & 三格距、混合重整化、联合外推、误差 & integrate；TMD、准 TMD 与软结构 \\
  U084 & \texttt{胶子PDF连续极限幻灯片.tex} & 工作流、数据布局、GPU/MPI 说明 & deduplicate；U076/\allowbreak U083/\allowbreak U086 \\
  U085 & \texttt{胶子PDF连续极限论文.tex} & 与 U083 相同的正文、公式和补充材料 & deduplicate；U083 \\
  U086 & \texttt{胶子TMD梯度流连续极限.tex} & flow-TMD 推导、流时窗口、候选工作流 & integrate；梯度流 TMD、匹配与连续极限 \\
  U087 & \texttt{质子自旋危机解析.tex} & EMC 历史、两套自旋分解、证据争议 & integrate；质子自旋、胶子极化与连续极限 \\
\end{longtable}

\section{合并前的最低检查单}

\begin{enumerate}[leftmargin=*]
  \item \textbf{算符几何：}Wilson 线路径、方向、长度、端点和色表示是否与裸量、减除量及匹配核一致。
  \item \textbf{Lorentz 投影：}是否包含迹项、对偶张量符号、\(z\)-奇组合和核子自旋投影；是否把诊断单项误当 twist-2 本征算符。
  \item \textbf{UV 方案：}是否报告短距参考点、强子矩阵元归一化、\(Z_R\) 或 Wilson 圈、格距依赖和方案转换。
  \item \textbf{软与 rapidity：}TMD 是否除以正确的 \(S_I\) 平方根、使用同一表示 CS 核，并验证 \(\zeta\) 方案一致性。
  \item \textbf{匹配：}硬核阶数、Fourier 归一化、味单态混合、匹配标度和高阶误差是否写明。
  \item \textbf{极限：}是否同时报告大 \(z\) 尾部、多个 \(P_z\)、至少三格距的联合外推；flow 研究是否扫描多个流时。
  \item \textbf{误差：}统计、拟合、重正化、匹配、flow、手征、体积和外推误差是否分别给出并可传播到最终协方差。
  \item \textbf{证据等级：}结论是否停在表 \ref{tab:b05-claim-levels} 的已闭合层级，是否把 quark 先例或流程图误称为 gluon TMD 数值结果。
\end{enumerate}

\section{本批次保留的未解决项}

\begin{enumerate}[leftmargin=*]
  \item gluon TMD 的软函数和 CS 核如何以伴随表示与 quark 方案统一，尚无 B05 内的直接数据闭合。
  \item 大 \(b_T\) 区固定阶 OPE 失效后，格点数据、非微扰软核和现象学模型的误差如何分离。
  \item 味单态胶子--夸克混合矩阵的非微扰确定及其与 NLO/NNLO 匹配核的一致性。
  \item 梯度流与 rapidity/soft 方案的联合定义；流时外推和匹配阶数的相关误差尚缺多流时数值表。
  \item \(M_{ti;it}+M_{ij;ji}\) 中 \(M_{gg}\) 与接触项的精确约定，以及不同格点 Clover 定向的归一化一致性。
  \item 自旋预算的跨方案、跨标度协方差；现有预算区间不能直接作为闭合的物理结论。
\end{enumerate}

本附录不把“缺少一项”解释为整个框架无效。相反，B05 的可靠结论是：胶子算符、TMD 减除、软核、连续外推和自旋分解已经形成可审计的方法骨架；只有逐层满足上表条件，骨架才能升级为物理分布或自旋分量。
